\setcounter{section}{-1}
\section{Introducción a la administración pública}

Tema de introducción a la administración pública y el marco normativo. La administración pública es la empresa en la que vamos a trabajar, y estaría bien empezar a conocerla.

\localtableofcontents

\subsection{Introducción}


\subsection{Administración pública del estado}

    La administración es el brazo ejecutor de los fines políticos de un gobierno.

    El Gobierno reparte sus responsabilidades en los ministerios, que cumplen con los objetivos políticos designados. En españa hay alrededor de 8.000 gobiernos, entre el general, autonómicos, provinciales y locales. La administración pública está subordinada al gobierno y tiene unos límites precisos, realiza la función administrativa y de gestión del estado y entes públicos de distinto ámbito.

    \textit{Las administraciones gestionan las políticas de los gobiernos}

    Una administración se compone de:
    \begin{itemize}
        \item Personas físicas
        \item Medios económicos
        \item Organización y ordenación racional de los medios
        \item Fines y principios
        \item Actuación (lícita dentro de una competencia de órgano actuante)
    \end{itemize}

    La administración siempre debe actuar al amparo de la ley y el derecho.

\subsubsection{Concepto de Persona y Personalidad}

    Desde el punto de vista jurídico, es una entidad con derechos y obligaciones que puede actuar jurídicamente.

    Una empresa es una personalidad jurídica cuando adquiere la condición de ser reclamador de sus derechos y de responder de sus deberes. Cuando está constituida y tiene una razón y denominación social. En el caso de una empresa, los derechos se ejercitan por los dueños o consejo administrativo o representantes. Y es el capital social el que responde de los deberes y los dueños de manera subsidiaria.

    Una \textit{personalidad jurídica plena} es aquella, en virtud de la cual, se obtiene la capacidad para ejercer derechos, realizar actividades o adquirir obligaciones, sin limitación geográfica o competencial.

\subsubsection{Art. 103}
    El Artículo 103 de la constitución se enmarca dentro del título IV (Del gobierno y la administración). Regula la administración pública, sus principios rectores, su creación y funcionamiento y acceso.

    La administración debe servir con los principios de:
    \begin{description}
        \item[Objetividad]
        \item[Eficacia] La capacidad de alcanzar el objetivo deseado tras realizar una acción
        \item[Jerarquía] La organización es piramidal y la subordinación es estricta.
        \item[Descentralización] Supone transferir el poder del gobierno central a las autoridades que no están subordinadas.
        \item[Desconcentración] Paso del ámbito estatal al autonómico y local de las competencias que, en cierto modo, si están subordinadas.
        \item[Coordinación] La ley 40/2015 del Régimen Jurídico del Sector Público Trata en su título 3° sobre la obigación de la coordinación de las administraciones. La ley obliga a las administraciones a trabajar entre ellas.
    \end{description}

    \noindent\fbox{\parbox{\textwidth}{Según los principios de actuación, ¿la descentralización implica que las competencias emanan del gobierno central? Por ejemplo, en el caso de los Mossos, es el estado el que le cede estas competencias a Cataluña? puede retomar unilateralmente estas competencias de vuelta? ¿Se hace de la misma manera o mediante el artículo 153?}}

    El RDL 5/2015 regula el estatuto de los funcionarios públicos. El acceso a la función pública se regula de acuerdo con mérito y capacidad.

\subsubsection{Administración Local}
    La constitución Española, en su artículo 140, dice que los municipios son responsables de sus actos, siempre que tengan la competencia. El artículo 141 dice que es una entidad local, con personalidad jurídica propia, determinada por la agrupación de municipios.

    En las autonomías monoprovinciales no hay diputación, pues estas funciones las asume la autonomía.

    \noindent\fbox{\parbox{\textwidth}{Ciertas comunidades se formaron por procedimiento histórico (Cataluña, País Vasco) y otras por ciudades limítrofes}}

    Los artículos 148 y 149 regulan las competencias transferibles y no transferibles. El artículo 150 dice que hay excepciones a las rigideces de ambos artículos. Estas excepciones son las que dan pie a muchas consultas al constitucional.

\subsubsection{Constitución y Órganos constitucionales}
    La constitución prevee ciertos órganos que no dependen del gobierno, ministerios o del estado:
    \begin{itemize}
        \item La Corona
        \item Cortes Generales
        \item El gobierno
        \item Consejo General del Poder Judicial
        \item Tribunal Constitucional
        \item Tribunal de cuentas
        \item Consejo de Estado
        \item Defensor del Pueblo
        \item Ministerio Fiscal
        \item Consejo Económico y Social
    \end{itemize}

    Todos son órganos constitucionales, pero no son iguales. Tampoco son administración pública, pero España no estaría constituida como lo está ahora sin estos órganos.

    El Consejo Económico y Social no está nombrado en la constitución, pero la ley 21/1991 le da especial importancia a sindicatos y otras entidades.

    \noindent\fbox{\parbox{\textwidth}{El ministerio fiscal no es un ministerio}}

\subsubsection{La administración general del estado}

    La administración estatal se divide en la Administración General y la Administración Institucional. A su vez, la Administración general se divide en la Administración Central (Ministerios), Territorial (Comunidades autónomas, competencias locales) y la General en el Exterior.

    En el derecho público, las normas jurídicas regulan su funcionamiento. En el derecho privado, es el conjunto de normas jurídicas el que regula las relaciones entre particulares
    % Neceita expandir

    Normalmente se respeta el nombre con el que la ciudadanía conoce a un ente, por encima de la denominación. Aunque (En teoría) la denominación debe ser efectiva.

    La administración estatal se compone de
    \begin{itemize}
        \item Admin General del estado
        \begin{itemize}
            \item Administración central
            \item Organización territorial
            \item Administración General del estado en el exterior
        \end{itemize}
        \item Administración institucional del estado
        \begin{itemize}
            \item Organismos públicos
            \begin{itemize}
                \item Organismos autónomos
                \item Entidades públicas empresariales
                \item Agencias estatales
            \end{itemize}
            \item Autoridades administrativas independientes
            \item Sociedades mercantiles
        \end{itemize}
    \end{itemize}

    \noindent\fbox{\parbox{\textwidth}{INVENTE es una aplicación de hacienda que concentra la información de todo el sector institucional, estatal, autonómico y local.}}

    \paragraph{Entidades de Derecho Público}
        Entidades por derecho propio. Su existencia la dicta una norma jurídica, normalmente un Real Decrto. Llevan acciones que podrían encargarse a una Dirección General. Tienen tesorería propia y autónomía. Existen unas 80. Esta forme les da la autonomía de operar a parte del estado. Ejemplos: CIS, INAP, CSdD

    \paragraph{Entidades públicas empresariales}
        Entidades con personalidad, patrimonio y autonomía propias. Determinadas acciones se realizan acorde al derecho público, y parte al derecho privado. Se financian principalmente por el ingreso de mercado. Ejemplos: FNMT, ADIF, Renfe, Enaire, ICO

    \paragraph{Agencias estatales}
        Entidades con personalidad, patrimonio y autonomía propia. Determinadas acciones las realizan acorde al derecho privado, con ciertas particularidades. Hay autonomía funcional y es responsable de los resultados establecidos en la ley. El estado no entra al detalle en gastos, mas allá del gasto en personal, que son obligados. Tiene potestad administrativa nen lo que le dice la ley. Ejemplos: ABOE, AEMET

    \paragraph{Autoridades administrativas independientes}
        Entidades de derecho público con personalidad jurídica propia y funciones de regulación o supervisión externo sobre sectores económicos o actividades determinadas. Su actuación requiere independencia funcional o administrativa del estado. Ejemplo: AEPD, CNMC, CNMV, CTBG, CSN, FROB.

    \paragraph{Sociedad Mercantil}
        Entidades participadas socialmente y públicamente en mas de un 50\% por el estado. Ejemplos: Museo del Prado, Correos, RTVE, Paradores, Lotería, Aena.

    \paragraph{Consorcio} Entidades creadas por varias administraciones públicas para el desarrollo de actividades de interés común. Ejemplos: Consorcio de transporte de la comunidad de Madrid (EMT, Metro, Cercanías)

    \paragraph{Fundaciones}
        Tienen que cumplir al menos uno de los requisitos de estar constituida, integrada patrimonialmente o correspondida en derechos por mas de un 50\% públicamente. Tienen cargo a los PGdE. Ejemplos: Ramón Areces, Giner de los Ríos, Pablo Iglesias

    \paragraph{Fondo sin personalidad jurídica}
        Masa presupuestaria carente de responsabilidad jurídica. Se usa \textit{ad hoc} para la gestión de un fin concreto, como fondos de ayuda a una catastrofe.

\subsubsection{Organización de la Administración General del Estado}
    La administración General se organiza en el Gobierno, Consejo de Ministros, Organos Ministeriales, Organización territorial y Órganos en el exterior, como (misiones diplomáticas, delegaciones, oficinas consulares).

    Los Órganos superiores son Ministros y Secretarios de Estado, estos son opcionales. Por debajo están los órganos directivos: Subsecretarios y Secretarios Generales, Secretarios generales técnicos, Directores Generales y subdirectores generales. Los órganos directivos tienen que ser funcionarios. Todos son altos cargos y se regulan por la ley 3/2015.

    En cada ministerio hay al frente un ministro, nombrado por el rey a propuesta del presidente del gobierno. los órganos se dividen en superiores y directivos. Obligatoriamente, cada ministerio debe tener una subsecretaría homónima. Por debajo hay una secretaría general técnica. Las demás secretarías de estado y generales no son obligatorias, pero si convenientes y frecuentes.

    Dentro de las competencias de cada ministerio, los gobiernos pueden hacer y deshacer, mover secretarías y funciones. Tanto dentro como entre ministerios. Los cargos políticos se deciden personalmente y con la misma categoría de ley, pueden cambiar. Según baja el escalafón los cargos son menos políticos.

    Esta información se encuentra en los portales de transparencia.

\subsection{Prelación de leyes}

    Es el conjunto de normas jurídicas que conforman el ordenamiento jurídico español.

    Se habla de ley hablando de normas jurídicas.

    Las fuentes del derecho por orden son: Ley escrita, Costumbre, Principios Generales del Derecho, y la Jurisprudéncia.

\subsubsection{Jerarquía normativa}

    No todas las leyes nacen iguales, ni se aprueban igual. Ciertas leyes tienen prioridad sobnre otras, según la dificultad de aprobarlas:

    \begin{itemize}
        \item Derecho comunitario
        \item Constitución
        \item Tratados internacionales
        \item Leyes y Leyes Autonómicas
        \item Real Decreto Ley y Real Decreto Administrativo
        \item Reglamentos
        \item Costumbre
        \item Principios GEnerales del Derecho
    \end{itemize}

    Los Reales Decretos y Reales Decreto Ley emanan del ejecutivo y tienen rango de ley. Los Reales Decretos Legislativos pueden ser textos propios o agrupación de normas. Los Reales Decreto Ley son normas que se aprueban sin ir al congreso. Las costumbres y PGdD son normas no escritas.

    Muchas normas relegan a normativa inferior los detalles de su aplicación. Permitiendo amoldarlos a la realidad sin tener que pasar por el proceso de cambiar la norma superior, mientras protege las líneas generales.

    \paragraph{Normas comunitarias}
        También llamado derecho comunitario. es el conjunto de normas y principios que determinan el funcionamiento de la unión europea. No es lo mismo que el derecho internacional.

        \subparagraph{Reglamento europeo}
            Norma jurídica que emana de las instituciones. Tiene efecto directo y está por encima de las normas nacionales. Tiene alcance general y se aplica sin ninguna transposición.

        \subparagraph{Directivas}
            Instrumento indirecto. Es una norma dirigida a los países para que legislen de una manera. La Directiva otorga un plazo para que se transponga la ley. Los estados deben interpretar y transponer. Estas leyes resultantes de las directivas terminan en las fronteras, mientras que los reglamentos se aplican en el conjunto de la unión europea.

        \subparagraph{Decisiones}
            Norma similar al reglamento, dirigida a instituciones u órganos concretos. Por ejemplo el reglamento IDAS

    \paragraph{Constitución española}
        Establece el marco jurídico y de competencia de las Comunidades Autónomas a la hora de dictar normas con rango de ley. En un trámite conjunto entre estado y autonomías se tramita los EStatutos de Autonomía, que tienen el rango de Ley Orgánica.

        La constitución indica el proceso de elaboración de leyes (Art. 81 a 92) y las formas de iniciativa legislativa (Arts 86 a 88). Si la ley la presenta el gobierno se considera un proyecto de ley, si no es una proposición de ley o una Iniciativa Legislativa Popular. Siendo preferente el Proyecto de Ley

    \paragraph{Tratado internacional}
        Acuerdo entre estados regulado en el derecho internacional. Requiere que el acuerdo incluya al menos dos entidades jurídicas internacionales, es decir, no solo estados, también organismos internacionales como la ONU. La Constitución prevalece sobre los tratados internacionales. Esta contradicción se puede requererir por parte del gobierno o cámaras al Tribunal Constitucional.

        Los tratados internacionales se publican como tal, no como ley, y forman parte del ordenamiento interno.

    \paragraph{Ley Orgánica}
        Requiere el voto favorable de la mayoría absoluta del congreso en votación final. Regulan determinadas materias, especialmente aquellas referidas a derechos fundamentales.
        
        Están a la altura de las leyes ordinarias. Se diferencian en su regulación. Quedan excluidas de los Reales Decretos Legislativos, Decretos Ley o Iniciativa Legislativa Popular.

    \paragraph{Ley Ordinaria}
        Se aprueba por mayoría sumple y recogen las materias que regulan  las leyes orgánicas. Se tramitan por iniciativa de la cámara, del gobierno, o popular. Las leyes de las comunidades autónomas también son leyes.

    \paragraph{Real Decreto Ley}
        Norma que entra en vigor inmediatamente por motivos de urgencia y necesidad. No pueden regular temas sujetos a ley orgánica.

        En su entrada en vigor, no son aprobadas por las cortes. Se deben ratificar en un plazo de 30 días. Hasta entonces se encuentran en vigor y su no aprobación no deroga retroactivamente. Tienen rango de Ley se debe aplicar como tal. Se le pueden incorporar enmiendas.

        Es normal que un RD toque muchos aspectos de multitud de leyes

    \paragraph{Real Decreto Legislativo}
        Norma jurídica con rango de ley. Emana del ejecutivo, en virtud de delegación expresa efectuada por el legislativo.

        Suele usarse para la compilación de diversas leyes afectando a un mismo tema, o para la redacción de una ley de bases.

        El consejo de estado filtra los Reales Decretos Legislativos

    \paragraph{Reglamento}
        Norma escrita con rango inferior a la ley, se dicta por la Administración Pública en virtud de su competencia. Hay reglamentos que no desarrollan leyes.

        No hay materias reservadas a los reglamentos. La administración está obligada a seguir los reglamentos, ya que está sujeta a la ley y al resto del ordenamiento jurídico.

        Reciben el nombre de la jerarquía normativa:
        \begin{description}
            \item[Real Decreto] Presidencia o Consejo de Ministros
            \item[Orden ministerial] Uno o Varios Ministerios
            \item[Circular, Resolución, Instrucción, Órdenes de servicio] Otros órganos
        \end{description}

        Algunos ejemplos son:
        \begin{itemize}
            \item RD 828/2023 por el que se nombra al presidente
            \item RD 829.2023 por el que se reestructuran los ministerios
            \item RD 203/2021 por el que se aprueba el reglamento de funcionamiento electrónico del Sector público
        \end{itemize}

        \subparagraph{Reglamento ejecutivo}
            O \textit{Secundum Legem} Desarrollan la aplicación de una ley. De esta manera la ley define las líneas generales, y los detalles se pueden detallar en un reglamento. También pueden cambiarse estos detalles, quedando la ley fija.

        \subparagraph{Reglamento independiente}
            O \textit{Praeter legem} Se dictan en ausencia de una ley. Surgen al margen de una habilitación legal previa (suplen la inexistencia de una ley). Ha quedado relegado prácticamente a organización interna de la administración. Estos reglamentos no pueden crear derechos o imponer obligaciones que no estén marcados en la ley.

        \subparagraph{Reglamento de necesidad}
            O \textit{Contra legem} se dictan con caracter excepcional contra las normas vigentes, incluso por órganos distintos a los que tienen la potestad reglamentaria, por motivos de emergencia o anormalidad transitorias. Tienen vigencia hasta que se supera la situación excepcional.
